\chapter{Changes in \texorpdfstring{\ps}{pystencils}}\label{r03:ch:pystencils}

\emph{How \ps{} vectorizes the kernel for RVV, and the three additions that were adopted:
software prefetch, a new instruction selection with scalar operands and fused multiply-add,
and row pointers with constants kept in registers.}

The fork is \href{\rPSurl/tree/\rPSsha}{\texttt{mirrored-pystencil2}}, branch
\texttt{rvv-support}, commit \rcommitps. The changes described here touch ten files
(\href{\rPSurl/compare/9a989e1625fbb0272876f550a5f7d3456fdaca25...\rPSsha}{\texttt{9a989e16..a81dc547}}:
807 lines added, 2 removed). Only code for the RVV target changes. The new extension points
do nothing for other targets, and the prefetch option is refused for them.

\section{How the kernel is vectorized}\label{r03:sec:ps-simd}

\ps{} vectorizes the innermost loop, which runs over $x$; in the fzyx layout consecutive $x$
are consecutive in memory (\cref{r03:sec:bg-blocks}). One vector is one LMUL-1 register of
256~bits: 8 lanes in f32, 4 in f64. Each vector iteration therefore updates 8 (or 4) cells at
once. Every memory access is a unit-stride load or store of part of one row of one slice; the
kernel needs no gather or scatter instructions. If the row length is not a multiple of the
lane count, a scalar loop handles the remaining cells; at 64 cells per row it has nothing to do.
\Cref{r03:fig:simd} shows the layout and one iteration.

\begin{figure}[htbp]
  \centering
  \resizebox{\linewidth}{!}{% The SIMD scheme of the generated sweep: fzyx field, x innermost and contiguous, one LMUL-1 vector
% register (VLEN 256) = 8 f32 / 4 f64 lanes; 64 interior cells = 8 vector iterations per row, plus the
% scalar remainder loop pystencils always emits (empty here). Counts per iteration: the f32 odd AA
% kernel of Listing r03:lst:aa-odd (19 loads, 23 stores, 23 prefetches) and table r03:tab:selector-insns (~150 FP ops).
\begin{tikzpicture}[every node/.style={font=\scriptsize}]
  \def\rthreeCW{0.2}
  % ---- (a) one row of one PDF slice ----------------------------------------------------
  \node[anchor=west,font=\footnotesize\bfseries] at (-0.1,1.25) {(a) one row of one PDF slice in memory};
  \foreach \i in {0,...,67} {
    \pgfmathsetmacro{\xx}{\i*\rthreeCW}
    \ifnum\i=0 \def\rthreeFC{psgray!35}\else\ifnum\i>64 \def\rthreeFC{psgray!35}\else
      \pgfmathtruncatemacro{\chunk}{mod(div(\i-1,8),2)}
      \ifnum\chunk=0 \def\rthreeFC{psblue!22}\else\def\rthreeFC{psblue!8}\fi\fi\fi
    \ifnum\i>65 \def\rthreeFC{white}\fi
    \fill[\rthreeFC] (\xx,0) rectangle ++(\rthreeCW,0.45);
    \draw[psgray!60,line width=0.2pt] (\xx,0) rectangle ++(\rthreeCW,0.45);
  }
  \foreach \i in {66,67} \fill[pattern=north east lines,pattern color=psgray!60] (\i*\rthreeCW,0) rectangle ++(\rthreeCW,0.45);
  \draw[decorate,decoration={brace,amplitude=3pt}] (0.2,0.55) -- node[above=2pt] {interior $x = 0 \ldots 63$} (13.0,0.55);
  \node[anchor=north] at (0.1,-0.02) {ghost};
  \node[anchor=north east] at (13.6,-0.02) {ghost, 2 pad};
  \foreach \k in {0,...,7} \node[text=psblue!80!black] at (\k*8*\rthreeCW+0.2+4*\rthreeCW,0.225) {\tiny it.~\the\numexpr\k+1\relax};
  \node[r03label,anchor=west] at (0,-0.72)
    {68 floats = 272~B per row (rows padded to 16~B); 19 such slices, one per direction, in the fzyx layout};

  % ---- (b) one vector register ---------------------------------------------------------------
  \node[anchor=west,font=\footnotesize\bfseries] at (-0.1,-1.25) {(b) one vector register, LMUL 1, VLEN 256: one \texttt{vle32.v} loads one shaded block of (a)};
  \foreach \l in {0,...,7} \node[r03lane,minimum width=8mm] at (0.8+\l*0.8,-1.9) {$x{+}\l$};
  \node[anchor=west] at (7.25,-1.9) {\texttt{vfloat32m1\_t}: 8 lanes of 32 bits (f32)};
  \foreach \l in {0,...,3} \node[r03lane,minimum width=16mm,fill=psteal!15,draw=psteal!80!black] at (1.2+\l*1.6,-2.55) {$x{+}\l$};
  \node[anchor=west] at (7.25,-2.55) {\texttt{vfloat64m1\_t}: 4 lanes of 64 bits (f64)};

  % ---- (c) one vector iteration ----------------------------------------------------------------
  \node[anchor=west,font=\footnotesize\bfseries] at (-0.1,-3.3) {(c) one iteration of the vector loop (8 cells, f32)};
  \node[r03code,text width=138mm,align=left,anchor=north west] at (0,-3.55) {%
for x in 0, 8, 16, ... (step 8):            // then a scalar remainder loop for the last (64 mod 8 = 0) cells\\
\ \ 23 prefetches   prefetch\_read/write(\&row\_k[x + 32])     // 128 B ahead, one per access\\
\ \ 19 loads        v\_i = vle32(\&row\_i[x])                  // one row pointer per slice and offset\\
\ \ $\approx$150 FP ops moments, equilibrium, relaxation; 64 of them fused (vfmacc, vfmsac, ...)\\
\ \ 23 stores       vse32(\&row\_k[x], ...)                   // 19 PDF slots + rho + 3 velocity components};

  % ---- (d) register grouping tested ------------------------------------------------------------
  \node[anchor=west,font=\footnotesize\bfseries] at (-0.1,-5.75) {(d) register grouping (LMUL), tried on the hand AA kernel (Section~\ref{r03:sec:lever-sched})};
  \node[r03ok,text width=40mm,font=\scriptsize] at (2.1,-6.5) {LMUL 1: one register,\\8 f32 lanes; \textbf{used}};
  \node[r03no,text width=40mm,font=\scriptsize] at (6.9,-6.5) {LMUL 2: two registers,\\16 lanes; 216 vs 163 cycles/cell};
  \node[r03no,text width=40mm,font=\scriptsize] at (11.7,-6.5) {LMUL 4: four registers,\\32 lanes; 226--229, spills};
\end{tikzpicture}
}
  \caption[How the generated kernel uses the vector registers]{How the generated kernel uses
    the vector registers: (a) one row of one PDF slice in memory and its 8 vector iterations,
    (b) one vector register, (c) the operations of one iteration of the f32 AA loop, (d) the
    register groupings tried in \cref{r03:sec:lever-sched}.}\label{r03:fig:simd}
\end{figure}

\paragraph{Where the changes sit.} \ps{} builds a kernel in stages
(\cref{r03:fig:pipeline}). It first represents the loop nest as a syntax tree, its
\term{intermediate representation} (IR). Transformations then rewrite that tree. Finally
the tree is \term{lowered}: \ps{} replaces generic operations by target-specific ones
(here RVV intrinsics) and prints C++. The \term{platform} class of a target (\code{RvvCpu} for
RVV) supplies the target-specific parts. The three changes of this chapter are two new
transformations and a rewritten instruction selection.

\begin{figure}[htbp]
  \centering
  \resizebox{\linewidth}{!}{% The pystencils kernel-creation pipeline (codegen/driver.py, DefaultKernelCreationDriver.__call__)
% for a vector CPU target, with the parts added or changed by this work highlighted (teal).
\begin{tikzpicture}[every node/.style={font=\scriptsize},node distance=3.2mm]
  \node[r03box,text width=46mm,font=\scriptsize] (a) {update of one cell\\(assignments from lbmpy)};
  \node[r03box,text width=46mm,font=\scriptsize,below=of a] (b) {loop nest over all cells; innermost loop ($x$)\\to be vectorized: 8 lanes (f32), 4 (f64)};
  \node[r03box,text width=46mm,font=\scriptsize,below=of b] (c) {create the loops: vector loop\\+ scalar loop for the remainder};
  \node[r03box,text width=46mm,font=\scriptsize,below=of c] (d) {general optimisations: fold constants,\\move loop-invariant code out of loops};
  \node[r03ok,text width=46mm,font=\scriptsize,below=of d] (e) {\texttt{prepare\_vector\_loops}\\(RVV: \texttt{HoistRowPointers})};
  \node[r03ok,text width=46mm,font=\scriptsize,below=of e] (f) {\texttt{insert\_prefetches}\\(only with \texttt{prefetch\_distance})};
  \node[r03ok,text width=46mm,font=\scriptsize,below=of f] (g) {instruction selection\\(\texttt{SelectIntrinsicsRvv})};
  \node[r03ok,text width=46mm,font=\scriptsize,below=of g] (h) {print C++ code; helper functions in\\\texttt{pystencils\_runtime/rvv.hpp}};
  \foreach \u/\v in {a/b,b/c,c/d,d/e,e/f,f/g,g/h} \draw[r03arrow] (\u) -- (\v);
  \draw[decorate,decoration={brace,amplitude=4pt,mirror}] ([xshift=-2mm]e.north west) -- node[left=5pt,align=right] {lowering\\(vector targets)} ([xshift=-2mm]g.south west);

  % what each new part does
  \node[r03ok,text width=66mm,font=\scriptsize,anchor=west] (e2) at ([xshift=9mm]e.east)
    {\texttt{row = \&ptr[rest]} before the $x$ loop; accesses become \texttt{row[x]}
     (\cref{r03:fig:rowptr}). Other platforms: no change.};
  \node[r03ok,text width=66mm,font=\scriptsize,anchor=west] (f2) at ([xshift=9mm]f.east)
    {one \texttt{prefetch\_read/write(\&row[x + D])} per distinct access, before the first memory
     statement (\cref{r03:fig:prefetch}). Other platforms refuse.};
  \node[r03ok,text width=66mm,font=\scriptsize,anchor=west] (g2) at ([xshift=9mm]g.east)
    {new: scalar copies of uniform values, .vf forms, fused multiply-add, constants through
     \texttt{in\_register} (\cref{r03:fig:selector})};
  \node[r03ok,text width=66mm,font=\scriptsize,anchor=west] (h2) at ([xshift=9mm]h.east)
    {new runtime functions: \texttt{prefetch\_read}, \texttt{prefetch\_write}
     (Zicbop HINTs), \texttt{in\_register}};
  \foreach \u in {e,f,g,h} \draw[psgray!70,dashed] (\u.east) -- (\u2.west);
\end{tikzpicture}
}
  \caption[The stages of pystencils and the additions]{The stages in which \ps{} builds a
    kernel for a vector CPU (\texttt{codegen/driver.py}). Teal: added or changed by the work in
    this report.}\label{r03:fig:pipeline}
\end{figure}

\section{Software prefetch}\label{r03:sec:ps-prefetch}

\paragraph{Option.} \code{cpu.vectorize.prefetch_distance}: the prefetch distance in bytes.
The default \code{None} means no prefetching. Defined in
\ghps{src/pystencils/codegen/config.py}; supported for \code{Target.RISCV_RVV} only.

\paragraph{Where it runs.} \ghps{src/pystencils/codegen/driver.py}: during lowering for vector
targets, if the option is set, the driver calls
\code{platform.insert_prefetches(kernel_ast, prefetch_distance)}. This happens after
\code{prepare_vector_loops} (\cref{r03:sec:ps-rowptr}) and before the instruction selection.

\paragraph{The transformation.} \code{InsertPrefetches}
(\ghps{src/pystencils/backend/transformations/insert_prefetches.py}, new):
\begin{itemize}
  \item In every innermost loop that contains unit-stride vector memory accesses, it inserts
        one prefetch per distinct address (base pointer plus offset), right before the first
        statement that accesses memory.
  \item The prefetch is a write prefetch if the loop writes that address, and a read prefetch
        otherwise. The prefetched address lies \code{distance_bytes} ahead of the accessed
        one.
  \item It skips accesses whose offset depends on a variable that is defined inside the loop
        at or after the insertion point, and accesses with a stride other than one.
\end{itemize}

\paragraph{Per platform.} \ghps{src/pystencils/backend/platforms/generic_cpu.py}: the base
class of vector platforms, \code{GenericVectorCpu}, raises an error
(\code{MaterializationError}) in \code{insert_prefetches}, so the option fails loudly on x86.
\ghps{src/pystencils/backend/platforms/rvv.py}: \code{RvvCpu.insert_prefetches} emits calls
to \code{pystencils::runtime::rvv::prefetch_read} and \code{prefetch_write}.

\paragraph{Runtime functions.} \ps{} ships a C++ header with helper functions for generated
code, \ghps{src/pystencils/include/pystencils_runtime/rvv.hpp}. If the compiler is told that
the core has Zicbop (\code{-march=..._zicbop}), the two functions use the compiler's
\code{__builtin_prefetch}. Otherwise they emit the instructions directly, \code{ori x0, rs, 1}
and \code{ori x0, rs, 3}, by inline assembly. These are the Zicbop prefetch encodings, which
are HINTs: cores without Zicbop ignore them (\cref{r03:sec:bg-rvv}). So no build flag is
needed. This matters because waLBerla's compiler flags do not contain \code{zicbop}, and the
X60's ISA string does not list it although the core executes the prefetches
(\cref{r03:sec:mem}). In the waLBerla build, the object file of the \code{explicit_cse_pf}
sweeps contains 65 such instructions: 42 in \code{StreamCollide} and 23 in \code{InitPdfs}.

\begin{figure}[htbp]
  \centering
  \resizebox{\linewidth}{!}{% Software prefetch along one row (f32): an iteration loads 8 floats = 32 B; the prefetch for the same
% access goes 128 B = 32 floats ahead, i.e. 4 iterations ahead; one prefetch per access and iteration,
% so each 64-B line receives two. (f64: 4 lanes = 32 B per iteration, 128 B = 16 doubles, also 4 ahead.)
\begin{tikzpicture}[every node/.style={font=\scriptsize}]
  \def\rthreeCW{0.2}
  % 48 floats = 3 cache lines of 16 floats
  \foreach \i in {0,...,47} {
    \pgfmathsetmacro{\xx}{\i*\rthreeCW}
    \ifnum\i<8 \def\rthreeFC{psblue!35}\else\ifnum\i>31 \ifnum\i<40 \def\rthreeFC{psorange!35}\else\def\rthreeFC{white}\fi\else\def\rthreeFC{white}\fi\fi
    \fill[\rthreeFC] (\xx,0) rectangle ++(\rthreeCW,0.45);
    \draw[psgray!60,line width=0.2pt] (\xx,0) rectangle ++(\rthreeCW,0.45);
  }
  % cache line boundaries
  \foreach \l in {0,1,2,3} \draw[psgray!90!black,line width=1pt] (\l*16*\rthreeCW,-0.1) -- ++(0,0.65);
  \foreach \l in {0,1,2} \node[r03label] at (\l*16*\rthreeCW+8*\rthreeCW,-1.15) {64-B cache line};
  \node[anchor=south] at (4*\rthreeCW,0.5) {\texttt{vle32} at $x$};
  \node[anchor=south] at (36*\rthreeCW,0.5) {\texttt{prefetch} at $x + 32$};
  \draw[r03arrow,psorange!90!black] (4*\rthreeCW,0.9) to[bend left=20] node[above] {same iteration, 128~B ahead} (36*\rthreeCW,0.9);
  % iterations
  \foreach \k in {0,...,5} \draw[psblue!60,decorate,decoration={brace,mirror,amplitude=2pt}]
    (\k*8*\rthreeCW+0.02,-0.55) -- node[below=2pt] {\tiny it.\,$k{+}\k$} (\k*8*\rthreeCW+8*\rthreeCW-0.02,-0.55);
  \node[anchor=west,text width=52mm] at (9.9,0.25)
    {The line holding $x + 32 \ldots x + 39$ is requested 4 iterations before \texttt{vle32} needs it.
     One prefetch per distinct access: 42 per iteration in the pull loop, 23 in the AA loop
     (\texttt{prefetch.r} for read-only, \texttt{prefetch.w} for written addresses).};
\end{tikzpicture}
}
  \caption[Software prefetch distance]{Software prefetch along one row (f32): 128~bytes ahead
    is four vector iterations ahead.}\label{r03:fig:prefetch}
\end{figure}

\section{Instruction selection: scalar operands and FMA}\label{r03:sec:ps-selector}

\term{Instruction selection} is the lowering step that maps each vector operation of the IR to
an RVV intrinsic. The original selection, \code{SelectIntrinsicsRvv} in
\ghps{src/pystencils/backend/platforms/rvv.py}, mapped every operation to the .vv form, and
kept every constant as a vector. It gained two rewrites, which are always on (described in the
class's documentation string):
\begin{description}
  \item[Scalar operands.] A value is \term{uniform} if it is the same in every lane: constants,
        broadcasts of a scalar, scalar variables, and $+$, $-$, $\times$, $\div$ and
        negation of uniform values. Every floating-point vector variable with a uniform value
        gets a scalar copy, declared right before it. An operation with one uniform operand then
        uses the .vf form with the scalar copy (\code{vfadd_vf}, \code{vfsub_vf},
        \code{vfmul_vf}, \code{vfdiv_vf}; \code{vfrsub_vf} or \code{vfrdiv_vf} when the scalar
        is the left operand). Variables that are assigned again later in the kernel, such as
        the running sum of a reduction, could change their value; a scan before the rewrite
        finds them and excludes them.
  \item[Fused multiply-add.] $a b + c$ becomes \code{vfmacc}, $a b - c$ becomes \code{vfmsac}
        (also when sympy writes it as $(-1) c + a b$), and $c - a b$ becomes \code{vfnmsac},
        each in the .vv or the .vf form. A uniform addend $c$ is not fused, because the
        accumulator of an FMA must be a vector, which would need a splat. The compiler does not
        form FMAs from intrinsics by itself (\cref{r03:sec:bg-rvv}), so the selection must.
\end{description}

\begin{figure}[htbp]
  \centering
  \resizebox{\linewidth}{!}{% The new RVV instruction selection (section 6.3) on one expression of the generated AA kernel
% (xi_23 = xi_22 + 1/8 * u0Mu1^2, Listing r03:lst:aa-odd), and what it does to the registers.
% Old mapping: one vv intrinsic per IR operation, constants as hoisted vector declarations.
\begin{tikzpicture}[every node/.style={font=\scriptsize}]
  \node[font=\footnotesize\bfseries] at (3.6,2.45) {before (one \texttt{vv} intrinsic per operation)};
  \node[font=\footnotesize\bfseries] at (11.4,2.45) {after (\texttt{.vf} operand, fused multiply-add)};
  \node[r03code,text width=68mm] (old) at (3.6,1.1) {%
c8 = vfdiv\_vv(vfmv(1.0f), vfmv(8.0f))   // hoisted\\
...\\
t1 = vfmul\_vv(u0Mu1, u0Mu1)\\
t2 = vfmul\_vv(t1, c8)              // c8 in a vector register\\
xi\_23 = vfadd\_vv(xi\_22, t2)};
  \node[r03code,text width=68mm] (new) at (11.4,1.1) {%
c8\_s = in\_register(1.0f / 8.0f)    // hoisted, FP register\\
...\\
t1 = vfmul\_vv(u0Mu1, u0Mu1)\\
xi\_23 = vfmacc\_vf(xi\_22, c8\_s, t1)  // +c8\_s*t1\\
\ };
  \draw[r03arrow] (old.east) -- (new.west);

  % register files
  \def\rthreeReg#1#2#3#4{% x, y, fill for the first #4 boxes, label
    \foreach \r in {0,...,31} {
      \pgfmathtruncatemacro{\rr}{\r}
      \ifnum\rr<#3 \def\rthreeFC{#4}\else\def\rthreeFC{white}\fi
      \fill[\rthreeFC] (#1+\r*0.21,#2) rectangle ++(0.19,0.32);
      \draw[psgray!60,line width=0.2pt] (#1+\r*0.21,#2) rectangle ++(0.19,0.32);
    }
  }
  \node[anchor=west] at (0,-0.25) {32 vector registers (orange: hoisted constants)};
  \rthreeReg{0}{-0.75}{31}{psorange!45}
  \node[r03label,anchor=north west,text width=68mm] at (0,-0.95)
    {31 constants compete with PDFs and temporaries for the 32 registers: 23 reloads + 7 spills
     per iteration (f32 AA + prefetch, even step)};
  \node[anchor=west] at (7.8,-0.25) {32 FP registers (teal: constants via \texttt{in\_register})};
  \rthreeReg{7.8}{-0.75}{9}{psteal!45}
  \node[r03label,anchor=north west,text width=68mm] at (7.8,-0.95)
    {scalar operands in the FP registers that a vector loop otherwise leaves idle (9 in the odd
     step of \cref{r03:lst:aa-odd}); vector spills and reloads 30 $\rightarrow$ 0, instructions
     471 $\rightarrow$ 341 per iteration (GCC~16)};
\end{tikzpicture}
}
  \caption[The new instruction selection]{The new instruction selection on one expression of
    the generated kernel (the old form is a schematic of the original mapping), and its effect
    on the registers.}\label{r03:fig:selector}
\end{figure}

\begin{table}[htbp]
  \centering\small
  \caption{Effect on the vector loops: instructions per iteration of 8 cells (f32) or 4 cells
    (f64), compiled with the GCC~16 cross compiler. Spills and reloads: vector registers saved to
    and loaded from the stack.}\label{r03:tab:selector-insns}
  \begin{tabular}{@{}lrrrr@{}}
    \toprule
    \textbf{Kernel} & \textbf{before} & \textbf{after} & \textbf{vector instr.} & \textbf{spills + reloads}\\
    \midrule
    f32 AA + prefetch, even step & 471 & 341 & 280 $\rightarrow$ 235 & 30 $\rightarrow$ 0\\
    f32 AA + prefetch, odd step & 556 & 396 & 284 $\rightarrow$ 236 & 34 $\rightarrow$ 0\\
    f64 AA + prefetch, even step & 473 & 350 & 280 $\rightarrow$ 235 & 30 $\rightarrow$ 0\\
    f32 pull + prefetch & 458 & 449 & 271 $\rightarrow$ 239 & 3 $\rightarrow$ 0\\
    \bottomrule
  \end{tabular}
\end{table}

The floating-point instructions per iteration drop from 207 to about 150, 64 of them FMAs.
Under QEMU the f64 checksums are unchanged to 13 digits. In f32 they change by about
$10^{-8}$ relative, because an FMA rounds once where a multiplication followed by an addition
rounds twice. On the board the f32 AA kernel with prefetch becomes 22\,\% faster
(\cref{r03:tab:selector}).

\section{Row pointers and constants in registers}\label{r03:sec:ps-rowptr}

\paragraph{\texttt{HoistRowPointers}}
(\ghps{src/pystencils/backend/transformations/hoist_row_pointers.py}, new). Consider a
unit-stride vector access \code{ptr[ctr + rest]} in an innermost loop, where \code{ctr} is the
loop counter (or a copy of it) and \code{rest} does not change inside the loop. The
transformation declares \code{row = &ptr[rest]} before the loop and rewrites the access to
\code{row[ctr]}. Accesses with the same pointer and the same \code{rest} share one row pointer.
The address computation then happens once per row instead of once per iteration
(\cref{r03:fig:rowptr}).

\begin{figure}[htbp]
  \centering
  \resizebox{\linewidth}{!}{% HoistRowPointers: before, every access computes its address from the base pointer inside the
% x loop, and GCC shares one pointer register between neighbours such as x-1 and x+1, adding the offset
% in every iteration (RVV loads/stores have no immediate offset); after, one row pointer per distinct
% (pointer, rest) is set once per row. Instruction counts: f32 AA odd step, GCC 16 cross.
\begin{tikzpicture}[every node/.style={font=\scriptsize}]
  \node[font=\footnotesize\bfseries] at (3.6,2.35) {before};
  \node[font=\footnotesize\bfseries] at (11.4,2.35) {after \texttt{HoistRowPointers}};
  \node[r03code,text width=72mm] (b) at (3.6,0.9) {%
for z, y:\\
\ \ for x in 0, 8, ...:\\
\ \ \ \ v3 = vle32(\&f[x - 1 + y*s1 + z*s2 + 3*s3])\\
\ \ \ \ v4 = vle32(\&f[x + 1 + y*s1 + z*s2 + 4*s3])\\
\ \ \ \ ...  // GCC: pointers based on x, +-4 B added\\
\ \ \ \ \ \ \ \ \ // per iteration (addi), extra addresses spilled};
  \node[r03code,text width=70mm] (a) at (11.4,0.9) {%
for z, y:\\
\ \ row3 = \&f[-1 + y*s1 + z*s2 + 3*s3]   // once per row\\
\ \ row4 = \&f[+1 + y*s1 + z*s2 + 4*s3]\\
\ \ for x in 0, 8, ...:\\
\ \ \ \ v3 = vle32(\&row3[x])    // pointer + x only\\
\ \ \ \ v4 = vle32(\&row4[x])};
  \draw[r03arrow] (b.east) -- (a.west);
  \node[r03no,text width=60mm] at (3.6,-0.75) {383 instructions per 8 cells (odd step)};
  \node[r03ok,text width=60mm] at (11.4,-0.75) {305 instructions per 8 cells (odd step); 23 row pointers};
\end{tikzpicture}
}
  \caption[Row pointers]{\texttt{HoistRowPointers}: how addresses are formed before and after
    (f32 AA odd step, GCC~16 cross compiler).}\label{r03:fig:rowptr}
\end{figure}

\paragraph{Extension point.} \code{GenericVectorCpu.prepare_vector_loops}
(\ghps{src/pystencils/backend/platforms/generic_cpu.py}) is a new method of the vector
platforms. The base class leaves the loops unchanged; \code{RvvCpu} applies
\code{HoistRowPointers}. The driver calls it before it inserts the prefetches, so the
prefetches use the row pointers too (\cref{r03:lst:aa-odd}).

\paragraph{Constants in registers.} Scalar copies (\cref{r03:sec:ps-selector}) whose value is a
constant are declared through \code{pystencils::runtime::rvv::in_register(x)}
(\ghps{src/pystencils/include/pystencils_runtime/rvv.hpp}). This function returns its argument
after an empty inline-assembly statement, \code{asm("" : "+f"(x))}, which tells the compiler
that the value may have changed. The compiler can then no longer treat it as a known constant,
and GCC keeps it in a floating-point register instead of loading it from the constant table for
every use.

\rthreeneedspace{12\baselineskip}
\lstinputlisting[style=cpp,caption={[Generated odd AA step with row pointers, prefetch, .vf and FMA]
  Excerpt of a generated kernel: the odd AA step of the D3Q19 update, f32, prefetch 128~bytes,
  from the lab file \texttt{misc/x60-lab/lab/src/f32-rvv-cse+aa+bpf=128+vf2+cg.cpp} (the file
  name is explained in \cref{r03:sec:repro-gen}). The two header comments and the shortened
  signature are added; all other lines are copied verbatim, and comments starting with
  ``\texttt{...}'' mark omitted lines.},
  label={r03:lst:aa-odd},nolol=false]{aa-odd-excerpt.cpp}

The kernel in \cref{r03:lst:aa-odd} has 23 row pointers (19 PDF slots, the density and three
velocity components), 23 prefetches and 19 loads per iteration, and nine constants that exist
both as an \code{in_register} scalar and as a vector.

The instruction counts drop (\cref{r03:sec:lever-rp}). In the lab on the board the cycles did
not; in the application the variants with prefetch became faster: AA by 2.5\,\% (f32) and
2.2\,\% (f64), pull by 5.2\,\% and 8.7\,\% (\cref{r03:tab:res-cg}). The change was kept
(\cref{r03:sec:res-cg}).

\section{Tests}\label{r03:sec:ps-tests}

\begin{table}[htbp]
  \centering\small
  \caption{Tests added or changed in \ps{} (directory \texttt{tests/nbackend/}).}\label{r03:tab:ps-tests}
  \begin{tabularx}{\linewidth}{@{}lX@{}}
    \toprule
    \textbf{Test} & \textbf{Checks}\\
    \midrule
    \multicolumn{2}{@{}l}{\ghps{tests/nbackend/test_rvv_codegen.py} (runs anywhere: inspects the
      generated code)}\\
    \code{test_prefetch_off_by_default} & no prefetch without the option\\
    \code{test_prefetch_one_per_distinct_access} & one prefetch per distinct access, read or write,
      distance in elements\\
    \code{test_prefetch_read_and_written_address_gets_write_prefetch} & write prefetch for an
      address that is both read and written\\
    \code{test_prefetch_needs_vector_platform} & the option is refused on other platforms\\
    \code{test_float_arithmetic} & updated for the new instruction selection\\
    \code{test_scalar_operands_use_vf_forms} & .vf forms for uniform operands\\
    \code{test_fused_multiply_add} & the FMA forms\\
    \code{test_scalar_addend_is_not_fused} & no FMA with a uniform addend\\
    \code{test_hoisted_constants_are_scalar_operands} & scalar copies of constants, through
      \code{in_register}\\
    \code{test_reassigned_vector_symbol_keeps_its_register} & no scalar copy for variables that are
      assigned again\\
    \code{test_stencil_accesses_use_row_pointers} & row pointers before the loop, with the stencil
      offsets moved into them\\
    \code{test_generated_code_compiles} & cross-compiles the kernels, now including one with
      prefetch\\
    \midrule
    \multicolumn{2}{@{}l}{\ghps{tests/nbackend/test_rvv_execution.py} (runs on RVV hardware only)}\\
    \code{test_prefetch_does_not_change_results} & identical results with and without prefetch, f32
      and f64\\
    \code{test_vf_forms_and_fma_match_numpy} & every form against numpy, within a tolerance in
      units of the last place (ulp)\\
    \code{test_stencil_row_pointers_match_generic_cpu} & a 3D stencil with row pointers against
      the plain C++ target\\
    \bottomrule
  \end{tabularx}
\end{table}

Two of these tests were checked by \term{mutation}: the code they protect was broken on purpose,
and the test had to fail. \code{test_reassigned_vector_symbol_keeps_its_register} fails when the
exclusion of re-assigned variables is removed, and \code{test_stencil_accesses_use_row_pointers}
fails when \code{prepare_vector_loops} is disabled. The test results are in
\cref{r03:sec:ver-ps}.
